\begin{lemma}
\label{lemma-one-orbit-geometric-galois}
Let $A$ be a normal domain with fraction field $K$.
Let $L/K$ be a (possibly infinite) Galois extension.
Let $G = \text{Gal}(L/K)$ and let
$B$ be the integral closure of $A$ in $L$.
\begin{enumerate}
\item For any two primes
$\mathfrak q, \mathfrak q' \subset B$ lying over the same prime in $A$
there exists a $\sigma \in G$ with $\sigma(\mathfrak q) = \mathfrak q'$.
\item Let $\mathfrak q \subset B$ be a prime lying over
$\mathfrak p \subset A$. Then $\kappa(\mathfrak q)/\kappa(\mathfrak p)$
is an algebraic normal extension and the map
$$
D = \{\sigma \in G \mid \sigma(\mathfrak q) = \mathfrak q\}
\longrightarrow
\text{Aut}(\kappa(\mathfrak q)/\kappa(\mathfrak p))
$$
is surjective.
\end{enumerate}
\end{lemma}

\begin{proof}
Proof of (1). Consider pairs $(M, \sigma)$ where $K \subset M \subset L$
is a subfield such that $M/K$ is Galois, $\sigma \in \text{Gal}(M/K)$
with $\sigma(\mathfrak q \cap M) = \mathfrak q' \cap M$.
We say $(M', \sigma') \geq (M, \sigma)$ if and only if
$M \subset M'$ and $\sigma'|_M = \sigma$.
Observe that $(K, \text{id}_K)$ is such a pair as $A = K \cap B$
since $A$ is a normal domain.
The collection of these pairs satisfies the hypotheses of Zorn's lemma,
hence there exists a maximal pair $(M, \sigma)$.
If $M \not = L$, then we can find
$M \subset M' \subset L$ with $M'/M$ nontrivial and finite and $M'/K$ Galois
(Fields, Lemma \ref{fields-lemma-normal-closure-inside-normal}).
Choose $\sigma' \in \text{Gal}(M'/K)$ whose restriction to $M$
is $\sigma$ (Fields, Lemma \ref{fields-lemma-galois-infinite}).
Then the primes $\sigma'(\mathfrak q \cap M')$ and $\mathfrak q' \cap M'$
restrict to the same prime of $B \cap M$. Since
$B \cap M = (B \cap M')^{\text{Gal}(M'/M)}$ we can
use Lemma \ref{lemma-one-orbit} to find $\tau \in \text{Gal}(M'/M)$
with $\tau(\sigma'(\mathfrak q \cap M')) = \mathfrak q' \cap M'$.
Hence $(M', \tau \circ \sigma') > (M, \sigma)$
contradicting the maximality of $(M, \sigma)$.

\medskip\noindent
Part (2) is proved in exactly the same manner as part (1). We
write out the details. Pick
$\overline{\sigma} \in \text{Aut}(\kappa(\mathfrak q)/\kappa(\mathfrak p))$.
Consider pairs $(M, \sigma)$ where $K \subset M \subset L$
is a subfield such that $M/K$ is Galois, $\sigma \in \text{Gal}(M/K)$
with $\sigma(\mathfrak q \cap M) = \mathfrak q \cap M$ and
$$
\xymatrix{
\kappa(\mathfrak q \cap M) \ar[r] \ar[d]_\sigma &
\kappa(\mathfrak q) \ar[d]_{\overline{\sigma}} \\
\kappa(\mathfrak q \cap M) \ar[r] & \kappa(\mathfrak q)
}
$$
commutes. We say $(M', \sigma') \geq (M, \sigma)$ if and only if
$M \subset M'$ and $\sigma'|_M = \sigma$.
As above $(K, \text{id}_K)$ is such a pair.
The collection of these pairs satisfies the hypotheses of Zorn's lemma,
hence there exists a maximal pair $(M, \sigma)$.
If $M \not = L$, then we can find
$M \subset M' \subset L$ with $M'/M$ finite and $M'/K$ Galois
(Fields, Lemma \ref{fields-lemma-normal-closure-inside-normal}).
Choose $\sigma' \in \text{Gal}(M'/K)$ whose restriction to $M$
is $\sigma$ (Fields, Lemma \ref{fields-lemma-galois-infinite}).
Then the primes $\sigma'(\mathfrak q \cap M')$ and $\mathfrak q \cap M'$
restrict to the same prime of $B \cap M$. Adjusting the choice
of $\sigma'$ as in the first paragraph, we may assume that
$\sigma'(\mathfrak q \cap M') = \mathfrak q \cap M'$.
Then $\sigma'$ and $\overline{\sigma}$ define maps
$\kappa(\mathfrak q \cap M') \to \kappa(\mathfrak q)$
which agree on $\kappa(\mathfrak q \cap M)$. Since
$B \cap M = (B \cap M')^{\text{Gal}(M'/M)}$ we can
use Lemma \ref{lemma-one-orbit-geometric}
to find $\tau \in \text{Gal}(M'/M)$ with
$\tau(\mathfrak q \cap M') = \mathfrak q \cap M'$
such that $\tau \circ \sigma$ and $\overline{\sigma}$
induce the same map on $\kappa(\mathfrak q \cap M')$.
There is a small detail here in that the lemma first
guarantees that $\kappa(\mathfrak q \cap M')/\kappa(\mathfrak q \cap M)$
is normal, which then tells us that the difference between
the maps is an automorphism of this extension
(Fields, Lemma \ref{fields-lemma-normal-embeddings-differ-by-aut}),
to which we can
apply the lemma to get $\tau$. Hence $(M', \tau \circ \sigma') > (M, \sigma)$
contradicting the maximality of $(M, \sigma)$.
\end{proof}